\clearpage\SourceLeaf{172}\pdfbookmark[0]{Methodus nova integralium valores per approximationem inveniendi}{seventh-paper}\thispagestyle{empty}
\begin{center}
METHODUS NOVA

\medskip
INTEGRALIUM VALORES

PER APPROXIMATIONEM INVENIENDI

AUCTORE

CAROLO FRIDERICO GAUSS

SOCIETATI REGIAE SCIENTIARUM EXHIBITA 1814. SEPT. 16.

\par\bigskip\rule{.7\linewidth}{.4pt}\par\medskip
Commentationes societatis regiae scientiarum Gottingensis recentiores. Vol.~\textsc{iii}.

Gottingae \textsc{mdcccxvi}.
\par\medskip\rule{.7\linewidth}{.4pt}
\end{center}
\clearpage\fancyhead[L]{\small METHODUS NOVA INTEGRALIUM VALORES}
\SourceLeaf{174}
\begin{center}
METHODUS NOVA

INTEGRALIUM VALORES

PER APPROXIMATIONEM INVENIENDI.

\medskip\rule{.1\linewidth}{.4pt}
\end{center}
\SeventhArticle{1}
Inter methodos ad determinationem numericam approximatam integralium
propositas insignem tenent locum regulae, quas praeeunte summo \textsc{Newton}
evolutas dedit \textsc{Cotes}. Scilicet si requiritur valor integralis $\int y\,\mathrm{d}x$
ab $x=g$ usque ad $x=h$ sumendus, valores ipsius $y$ pro his valoribus extremis
ipsius $x$ et pro quotcunque aliis intermediis a primo ad ultimum incrementis
aequalibus progredientibus, multiplicandi sunt per certos co\"efficientes
numericos, quo facto productorum aggregatum in $h-g$ ductum integrale
quaesitum suppeditabit, eo maiore praecisione, quo plures termini in hac
operatione adhibentur. Quum principia huius methodi, quae a geometris rarius
quam par est in usum vocari videtur, nusquam quod sciam plenius explicata
sint, pauca de his praemittere ab instituto nostro haud alienum erit.

\SeventhArticle{2}
Sit $n+1$ multitudo terminorum, quos in usum vocare placuit, statuamusque
$h-g=\Delta$, ita ut valores ipsius $x$ sint $g$, $g+\frac{\Delta}{n}$,
$g+\frac{2\Delta}{n}$, $g+\frac{3\Delta}{n}$ etc. usque ad $g+\Delta$,
respondeantque iisdem resp. valores ipsius $y$ hi $A$, $A'$, $A''$, $A'''$
etc. usque ad $A^{(n)}$: denique ponatur indefinite $x=g+\Delta t$, ita ut
$y$ etiam spectari possit tamquam functio ipsius $t$. Designemus per $Y$
functionem sequentem
\SourceLeaf{175}
\[
\begin{aligned}
&A.\;\frac{(nt-1)(nt-2)(nt-3)\ldotp\ldotp\ldotp\ldotp(nt-n)}{(-1)\quad(-2)\, .\,(-3)\ldotp\ldotp\ldotp\ldotp(-n)}\\
{}+{}&A'.\;\frac{nt\, .\,(nt-2)\, .\,(nt-3)\ldotp\ldotp\ldotp\ldotp\ldotp\ldotp(nt-n)}{1\, .\,(-1)\quad(-2)\ldotp\ldotp\ldotp\ldotp\ldotp\ldotp(1-n)}\\
{}+{}&A''.\;\frac{nt\, .\,(nt-1)\, .\,(nt-3)\ldotp\ldotp\ldotp\ldotp\ldotp\ldotp(nt-n)}{2.\quad1.\quad(-1)\ldotp\ldotp\ldotp\ldotp\ldotp\ldotp(2-n)}\\
{}+{}&A'''.\;\frac{nt\, .\,(nt-1)\, .\,(nt-2)\ldotp\ldotp\ldotp\ldotp\ldotp\ldotp(nt-n)}{3.\quad2.\quad1\ldotp\ldotp\ldotp\ldotp\ldotp\ldotp(3-n)}\\
{}+{}&\text{etc.}\\
{}+{}&A^{(n)}.\;\frac{nt(nt-1)\, .\,(nt-2)\ldotp\ldotp\ldotp\ldotp\ldotp\ldotp(nt-n+1)}{n\, .\,(n-1)\quad(n-2)\ldotp\ldotp\ldotp\ldotp\ldotp\ldotp1}
\end{aligned}
\]
sive $\sum\dfrac{A^{(\mu)}T^{(\mu)}}{M^{(\mu)}}$, ubi repraesentante $\mu$ singulos integros $0,1,2,3\ldotp\ldotp\ldotp n$,
\[
T^{(\mu)}=\frac{nt(nt-1)(nt-2)(nt-3)\ldotp\ldotp\ldotp\ldotp(nt-n)}{nt-\mu}
\]
\[
M^{(\mu)}\quad\text{valor ipsius }T\text{ pro }nt=\mu.
\]
Manifestum erit, $Y$ exhibere functionem algebraicam integram ipsius $t$ ordinis $n$, atque eius valores pro singulis $n+1$ valoribus ipsius $t$, puta $0,\frac1n,\frac2n,\frac3n\ldotp\ldotp\ldotp\ldotp1$ aequales esse valoribus ipsius $y$. Porro patet, si $Y'$ sit functio alia integra pro iisdem valoribus cum $y$ conspirans, $Y'-Y$ pro iisdem evanescere, adeoque per factores $t,t-\frac1n,t-\frac2n,t-\frac3n\ldotp\ldotp\ldotp\ldotp t-1$ et proin etiam per eorum productum (quod est ordinis $n+1$) divisibilem esse, unde patet, $Y'$, nisi prorsus identica sit cum $Y$, certo ad altiorem ordinem ascendere debere, sive $Y$ ex omnibus functionibus integris ordinem $n$ haud egredientibus unicam esse, quae pro illis $n+1$ valoribus cum $y$ conspiret. Quodsi itaque $y$, in seriem secundum potestates ipsius $t$ progredientem evoluta, ante terminum qui implicat $t^{n+1}$ omnino abrumpitur, cum $Y$ identica erit: si vero saltem tam cito convergit, ut terminos sequentes spernere liceat, functio $Y$ inter limites $t=0,t=1$ sive $x=g,x=h$ ipsius $y$ vice fungi poterit.

\SeventhArticle{3}
Iam integrale nostrum $\int y\,\mathrm{d}x$ transit in $\Delta\int y\,\mathrm{d}t$ a $t=0$ usque ad $t=1$ sumendum, cuius loco per ea, quae modo monuimus, adoptabimus $\Delta\int Y\,\mathrm{d}t$. Evolvendo itaque $T^{(\mu)}$ in
\[
\alpha t^n+\beta t^{n-1}+\gamma t^{n-2}+\delta t^{n-3}+\text{etc.}
\]
erit $\int T^{(\mu)}\,\mathrm{d}t$, a $t=0$ usque ad $t=1$,
\SourceLeaf{176}
\[
=\frac{\alpha}{n+1}+\frac{\beta}{n}+\frac{\gamma}{n-1}+\frac{\delta}{n-2}+\text{etc.}
\]
qua quantitate posita $=M^{(\mu)}R^{(\mu)}$, erit integrale quaesitum
\[
=\Delta(AR+A'R'+A''R''+A'''R'''+\text{etc.}+A^{(n)}R^{(n)})
\]
Exempli caussa apponemus computum co\"efficientis $R''$ pro $n=5$. Fit hic
\[
\begin{aligned}
T''&=5^5t^5-13\, .\,5^4t^4+59\, .\,5^3t^3-107\, .\,5^2tt+60\, .\,5\, .\,t\\[5pt]
M''&=2\times1\times(-1)\times(-2)\times(-3)=-12
\end{aligned}
\]
Hinc $-12R''=\allowbreak \frac{3125}{6}-1625+\allowbreak \frac{7375}{4}-\frac{2675}{3}+\allowbreak 150=\allowbreak -\frac{25}{12}$, adeoque $R''=\frac{25}{144}$.

Computus aliquanto brevior evadit, statuendo $2t-1=u$. Tunc fit
\[
T^{(\mu)}=\frac{\begin{gathered}(nu+n)(nu+n-2)(nu+n-4)\ldotp\ldotp\ldotp\\(nu-n+4)(nu-n+2)(nu-n)\end{gathered}}{2^n(nu+n-2\mu)}
\]
Ponamus
\[
\frac{\begin{gathered}(nnuu-nn)\, .\,(nnuu-(n-2)^2)\, .\,(nnuu-(n-4)^2)\\\, .\,(nnuu-(n-6)^2)\ldotp\ldotp\ldotp\ldotp\end{gathered}}{nnuu-(n-2\mu)^2}=U^{(\mu)}
\]
ubi numerator desinere debet in $\ldotp\ldotp\ldotp(nnuu-9)(nnuu-1)$, si $n$ est impar, vel in $\ldotp\ldotp\ldotp(nnuu-4)nu$, si $n$ est par, eritque
\[
T^{(\mu)}=\frac{(nu-n+2\mu)U^{(\mu)}}{2^n}
\]
Iam integrale $\int T^{(\mu)}\,\mathrm{d}t$ a $t=0$ usque ad $t=1$ acceptum aequale est integrali
\[
\int\tfrac12 T^{(\mu)}\,\mathrm{d}u=\int\frac{nuU^{(\mu)}\,\mathrm{d}u}{2^{n+1}}+\int\frac{(2\mu-n)U^{(\mu)}\,\mathrm{d}u}{2^{n+1}}
\]
ab $u=-1$ usque ad $u=+1$.

Statuendo itaque
\[
U^{(\mu)}=\alpha u^{n-1}+\beta u^{n-3}+\gamma u^{n-5}+\delta u^{n-7}+\text{etc.}
\]
\begingroup\rightskip=0pt plus1fil\relax\noindent (sponte enim patet, potestates $u^{n-2},\allowbreak u^{n-4},\allowbreak u^{n-6}$ etc. abesse), integralis pars $\int\dfrac{nuU^{(n)}\,\mathrm{d}u}{2^{n+1}}$ evanescet pro valore impari ipsius $n$, pars altera $\int\dfrac{(2\mu-n)U^{(\mu)}\,\mathrm{d}u}{2^{n+1}}$ vero pro valore pari, unde integrale $\int T^{(\mu)}\,\mathrm{d}t$ fiet pro $n$ pari\par\endgroup
\[
=\frac{n}{2^n}\left(\frac{\alpha}{n+1}+\frac{\beta}{n-1}+\frac{\gamma}{n-3}+\frac{\delta}{n-5}+\text{etc.}\right)
\]
\SourceLeaf{177}
pro $n$ impari autem
\[
=\frac{2\mu-n}{2^n}\left(\frac{\alpha}{n}+\frac{\beta}{n-2}+\frac{\gamma}{n-4}+\frac{\delta}{n-6}+\text{etc.}\right)
\]
In exemplo nostro habetur
\[
U''=(25uu-25)(25uu-9)=625u^4-850uu+225,\quad\text{adeoque}
\]
\[
-12R''=-\tfrac1{32}\left(125-\tfrac{850}{3}+225\right)=-\tfrac{25}{12}\quad\text{ut supra.}
\]
Observare convenit, fieri $U^{(n-\mu)}=U^{(\mu)}$, adeoque $\int T^{(n-\mu)}\,\mathrm{d}t=\pm\int T^{(\mu)}\,\mathrm{d}t$. signo superiore valente pro $n$ pari, inferiore pro impari. Quare quum facile perspiciatur, perinde haberi $M^{(n-\mu)}=\pm M^{(\mu)}$, semper erit $R^{(n-\mu)}=R^{(\mu)}$, sive e co\"efficientibus $R,R',R''\ldotp\ldotp\ldotp\ldotp R^{(n)}$ ultimus primo aequalis, penultimus secundo et sic porro.

\SeventhArticle{4}
Valores numericos horum co\"efficientium a \textsc{Cotesio} usque ad $n=10$ computatos ex \textit{Harmonia Mensurarum} huc adscribimus.

Pro $n=1$ sive terminis duobus.
\[
R=R'=\tfrac12
\]
Pro $n=2$ sive terminis tribus.
\[
R=R''=\tfrac16,\ R'=\tfrac23
\]
Pro $n=3$ sive terminis quatuor.
\[
R=R'''=\tfrac18,\ R'=R''=\tfrac38
\]
Pro $n=4$ sive terminis quinque.
\[
\begin{gathered}
R=R''''=\tfrac7{90},\quad R'=R'''=\tfrac{16}{45},\\[5pt]
R''=\tfrac2{15}
\end{gathered}
\]
Pro $n=5$ sive terminis sex.
\[
\begin{gathered}
R=R^{\mathrm{V}}=\tfrac{19}{288},\quad R'=R''''=\tfrac{25}{96},\\[5pt]
R''=R'''=\tfrac{25}{144}
\end{gathered}
\]
Pro $n=6$ sive terminis septem.
\[
\begin{gathered}
R=R^{\mathrm{VI}}=\tfrac{41}{840},\quad R'=R^{\mathrm{V}}=\tfrac9{35},\\[5pt]
R''=R^{\mathrm{IV}}=\tfrac9{280},\quad R'''=\tfrac{34}{105}
\end{gathered}
\]
Pro $n=7$ sive terminis octo.
\[
\begin{gathered}
R=R^{\mathrm{VII}}=\tfrac{751}{17280},\quad R'=R^{\mathrm{VI}}=\tfrac{3577}{17280},\\[5pt]
R''=R^{\mathrm{V}}=\tfrac{49}{640},\quad R'''=R''''=\tfrac{2989}{17280}
\end{gathered}
\]
Pro $n=8$ sive terminis novem.
\[
\begin{gathered}
R=R^{\mathrm{VIII}}=\tfrac{989}{28350},\quad R'=R^{\mathrm{VII}}=\tfrac{2944}{14175},\\[5pt]
R''=R^{\mathrm{VI}}=-\tfrac{464}{14175},\quad R'''=R^{\mathrm{V}}=\tfrac{5248}{14175},\\[5pt]
R^{\mathrm{IV}}=-\tfrac{454}{2835}
\end{gathered}
\]
\SourceLeaf{178}
Pro $n=9$ sive terminis decem.
\[
\begin{gathered}
R=R^{\mathrm{IX}}=\frac{2857}{89600},\quad R'=R^{\mathrm{VIII}}=\frac{15741}{89600},\\[5pt]
R''=R^{\mathrm{VII}}=\frac{27}{2240},\quad R'''=R^{\mathrm{VI}}=\frac{1209}{5600},\\[5pt]
R^{\mathrm{IV}}=R^{\mathrm{V}}=\frac{2889}{44800}
\end{gathered}
\]
Pro $n=10$ sive terminis undecim.
\[
\begin{gathered}
R=R^{\mathrm{X}}=\frac{16067}{598752},\quad R'=R^{\mathrm{IX}}=\frac{26575}{149688}.\quad
R''=R^{\mathrm{VIII}}=-\frac{16175}{199584},\\[5pt]
R'''=R^{\mathrm{VII}}=\frac{5675}{12474},\quad R^{\mathrm{IV}}=R^{\mathrm{VI}}=-\frac{4825}{11088},\\[5pt]
R^{\mathrm{V}}=\frac{17807}{24948}.
\end{gathered}
\]

\SeventhArticle{5}

Quum formula $\Delta(AR+\allowbreak A'R'+\allowbreak A'R'+\allowbreak A''R''+\allowbreak \text{etc.}+\allowbreak A^{(n)}R^{(n)})$ integrale
$\int y\,\mathrm{d}x$ ab $x=g$ usque ad $x=g+\Delta$, sive integrale
$\Delta\int y\,\mathrm{d}t$ a $t=0$ usque ad $t=1$ exacte quidem exhibeat,
quoties $y$ in seriem evoluta potestatem $t^n$ non transscendit, sed approximate
tantum, quoties $y$ ultra progreditur, superest, ut errorem, quem inducunt
termini proxime sequentes, assignare doceamus. Designemus generaliter per
$k^{(m)}$ differentiam inter valorem verum integralis $\int t^m\,\mathrm{d}t$
a $t=0$ usque ad $t=1$, atque valorem ex formula prodeuntem, ita ut sit
\[
\begin{aligned}
k&=1-R-R'-R''-R'''-\text{etc.}-R^{(n)}\\
k'&=\frac12-\frac1n(R'+2R''+3R'''+\text{etc.}+nR^{(n)})\\
k''&=\frac13-\frac1{nn}(R'+4R''+9R'''+\text{etc.}+nnR^{(n)})\\
k'''&=\frac14-\frac1{n^2}(R'+8R''+27R'''+\text{etc.}\ n^3R^{(n)})
\end{aligned}
\]
etc. Patet igitur, si $y$ evolvatur in seriem
\[
K+K't+K''tt+K'''t^3+\text{etc.}
\]
differentiam inter valorem verum integralis $\int y\,\mathrm{d}t$ atque valorem
approximatum formulae exprimi per
\[
Kk+K'k'+K''k''+K'''k'''+\text{etc.}
\]
Sed manifesto $k$, $k'$, $k''$ etc. usque ad $k^{(n)}$ sponte fiunt $=0$:
correctio itaque formulae approximatae erit
\[
K^{(n+1)}k^{(n+1)}+K^{(n+2)}k^{(n+2)}+K^{(n+3)}k^{(n+3)}+\text{etc.}
\]

Indolem quantitatum $k^{(n+1}$, $k^{(n+2)}$ etc. infra accuratius
perscrutabimur; hic sufficiat, valores numericos primae aut secundae, pro
singulis valoribus ipsius $n$, apposuisse, ut gradus praecisionis, quam formula
approximata affert, inde aestimari possit.
